% SPDX-License-Identifier: Apache-2.0
\section{A Unit That Answers the Reset}
\label{sec:x-ownrst}

A unit that only holds registers needs to know nothing about the
reset: while it is high, its registers go back to where they started.
A unit that wants to do something under reset, rather than merely
forget, takes \code{rst} as an \code{In<Bit>} and reads it, and then
what it does under reset is what its body says and nothing else. The
netlist gives it no clearing branch of its own, and the run does not
put its registers back on top of what the body drives. The lowering
tells the run which units these are: a unit whose \code{run} names a
port \code{rst} says so as it starts, and the run's reset then leaves
that unit's own registers to its body. A child that declares no
\code{rst} keeps the reset, as its module does (issue 878).

\code{Fetch} is the shape of a core's program counter. It counts up,
and under reset it loads its reset vector, \code{0x80}, rather than
going to zero. The run counts three, holds the reset for an edge, and
counts on from the vector. The port and the reset the netlist gives
every module are one net, so the run drives both at once.

\begin{figure}[htbp]
\centering
\input{ownrst_timing.tex}
\caption{The counter through a reset: \code{pc} counts, loads
\code{0x80} at the edge \code{rst} is high, and counts on from there.}
\label{fig:x-ownrst}
\end{figure}

\lstinputlisting[caption={\code{ex\_ownrst.rs}. Run with \code{bazel run //lib/examples:ex\_ownrst}.},label={lst:x-ownrst}]{lib/examples/ex_ownrst.rs}

\lstinputlisting[style=ascii,caption={What \code{ex\_ownrst} printed when the build ran it: the counter before, in and after the reset, then the netlist, with no reset branch of its own.},label={lst:x-ownrst-out}]{out_ownrst.txt}
